% !TeX root = ../../packet-and-principal-models.tex
\section{A common principal lattice and family}
The construction is a change of integral lattice, not a declaration that
the original trace polarization is principal. The subgroup at $7$ is
nonfree; requiring a common free line would incorrectly rule it out.
The special-fiber K-actions need not extend to a single global K-action.

\begin{proposition}[A common marked principal lattice at the three vertices]\label{prop:commonprincipal}
Use the common marking $\Lambda=\OL^2$, with $(a,b)\mapsto az_X+b$ at
$X=A,B,C$. There is one self-dual overlattice
$\Lambda\subset\Lambda^{\mathrm{common}}\subset\Lambda^\#$ of index
$453789$ stable under $\OL$ and all three specialized $\mathcal O_K$-actions
$\Sigma_X=S_X^{\mathsf T}$. With the respective complex structures it gives
three principally polarized isogenous special fibers with unchanged rational
CM Hodge structures.
\end{proposition}
\begin{proof}
Put $R=\OL/\mathfrak p_7^5=\mathbb F_7[u]/(u^5)$, where $u=t-1$.
The actual homology residues are
\[
 \Sigma_A\equiv\begin{pmatrix}1&3\\1&6\end{pmatrix},\qquad
 \Sigma_B\equiv\begin{pmatrix}6&2\\5&1\end{pmatrix},\qquad
 \Sigma_C\equiv\begin{pmatrix}5&5\\0&2\end{pmatrix}\pmod u.
\]
They all preserve the line of $v=(3,1)$. Define
\[
 H_7=u^2Rv+u^3R^2.
\]
Taking $w=(1,0)$, the vectors $v,w$ are an $R$-basis and
$H_7=u^2Rv\oplus u^3Rw$ has $\mathbb F_7$-basis
\[
 u^2v,\ u^3v,\ u^4v,\ u^3w,\ u^4w.
\]
Thus $|H_7|=7^5$.
For each $X$, write $\Sigma_Xv=\lambda_Xv+uq_X$; this immediately proves
$\Sigma_XH_7\subset H_7$. The determinant of any two elements of $H_7$
vanishes modulo $u^5$.

At $\mathfrak p_3$, all three $\Sigma_X$ vanish. Choose the same
first-coordinate $\mathbb F_{27}$-line $H_3$, and put $H=H_3\oplus H_7$.
To verify isotropy for the actual discriminant pairing, let $d=f'(t)$.
Multiplication by $d$ identifies $\Lambda^\#/\Lambda$ with
$(\OL/(d))^2$. Representatives $x,y\in\Lambda$ of elements of $H$ have
$\det(x,y)\in(d)$, by the two local determinant calculations. Hence
\[
 E_1(d^{-1}x,d^{-1}y)
 =-\operatorname{Tr}_{L/\Ql}\!\left(\frac{\det(x,y)}{d^2}\right)
 \in\mathbb Z,
\]
since the trace argument lies in the codifferent $d^{-1}\OL$.
As $|H|=3^3 7^5=\sqrt{|\Lambda^\#/\Lambda|}$, H is maximal isotropic.
Exact power-basis reduction verifies that all three
$U_X=(\Sigma_X-I)/2$ are integral. They preserve $\Lambda^\#$ as
$L$-linear integral matrices; since 2 is invertible on the discriminant
module, they preserve $H$. Thus its inverse image is simultaneously
$\OL$- and $\mathcal O_K$-stable and self-dual. The unchanged positive
Riemann form supplies the asserted principal polarizations.
\end{proof}

\begin{remark}[Free lines and the scope of the common construction]\label{rem:ABprincipallevels}
The common $H_7$ above is nonfree. Exhaustion of $\mathbb P^1(\mathbb F_7)$
shows that each pair of the three homology operators has only the common
residue line $(3,1)$. Its free lifts modulo $u^2$ have the form
$(3+au,1)$, with respective unique coefficients
\[
 a_A=3,\qquad a_B=6,\qquad a_C=2.
\]
Consequently no pair admits a common invariant \emph{free rank-one summand}
on $R^2$. This restricted obstruction does not exclude the common nonfree
maximal isotropic subgroup. Proposition~\ref{prop:commonprincipal} concerns
three specialized $K$-actions on one marked lattice. The global construction
below gives a family with RM and these three special fibers, without
asserting a single global $K$-action or a realization of Mostaed's literal
principally polarized locus.
\end{remark}

\begin{proposition}[Ambient and curve stabilizers of the common level]\label{prop:principalstabilizer}
Let $\Gamma=\mathrm{SL}_2(\OL)$, let $j:\Delta_0\to\Gamma$ be the
normalized integral representation of McMullen's quadrilateral group, and
let $H$ be the subgroup in Proposition~\ref{prop:commonprincipal}, with
$H_3=\mathbb F_{27}(1,0)$. For the transition representation
$\rho(g)=g^{-\mathsf T}$ put
\[
 \Gamma_H=\{g\in\Gamma:\rho(g)H=H\},
 \qquad \Delta_H=j^{-1}(\Gamma_H).
\]
Then
\[
 [\Gamma:\Gamma_H]=224,\qquad
 [\Delta_0:\Delta_H]=[\bar\Delta_0:\bar\Delta_H]=28.
\]
The exact ideal congruence level of $\Gamma_H$ is
$\mathfrak p_3\mathfrak p_7$: it contains $\Gamma(\mathfrak p_3\mathfrak p_7)$,
and $\Gamma(J)\subset\Gamma_H$ implies $J\subset\mathfrak p_3\mathfrak p_7$.
The number $28$ is the covering degree over the explicit parametrizing
orbifold $\mathbb H/\bar\Delta_0$, before an auxiliary neat refinement.
\end{proposition}
\begin{proof}
With the marking $(a,b)\mapsto az+b$, the isomorphism from the fiber at $z$
to the fiber at $gz$ is $w\mapsto w/(cz+d)$. Its homology matrix is
$g^{-\mathsf T}$. This transition action is a representation; it differs
from the fixed-point endomorphism matrix $g^{\mathsf T}$.

For any residue line $\ell\subset\mathbb F_7^2$, write
$H(\ell)=u^2Rv+u^3R^2$, where $v$ lifts $\ell$. An invertible matrix $T$
over $R$ preserves $H(\ell)$ precisely when its reduction preserves $\ell$.
Necessity follows by projecting to $u^2R^2/u^3R^2$. Conversely,
$Tv=\lambda v+uq$ gives $TH(\ell)\subset H(\ell)$, and finiteness gives
equality. Thus $\Gamma_H$ is the inverse image of two line stabilizers in
$\mathrm{SL}_2(\mathbb F_{27})\times\mathrm{SL}_2(\mathbb F_7)$.
Reduction is surjective: the finite groups are generated by elementary
matrices, whose entries lift by the Chinese remainder theorem. Hence the
ambient index is $(27+1)(7+1)=224$.

This description proves containment of $\Gamma(\mathfrak p_3\mathfrak p_7)$.
If $J$ is not contained in one of these primes, an elementary unipotent
with entry in $J$ can move the corresponding residue line. At $7$ use the
integral determinant-one basis with columns $(3,1),(-1,0)$. Applying
$\rho^{-1}$ gives an element of $\Gamma(J)$ outside $\Gamma_H$, proving
the asserted exactness of the level.

McMullen's generators $A,B,C,D$ generate $\Delta_0$ and satisfy $ABCD=I$
\cite[Sec.~4]{McMullen2026}, so $A,B,C$ suffice. Exact reduction of
$\rho(j(A)),\rho(j(B)),\rho(j(C))$ on the $224$-element set
$\mathbb P^1(\mathbb F_{27})\times\mathbb P^1(\mathbb F_7)$ fixes
$(3,1)$ in the second factor. Exhaustive closure of
$(\mathbb F_{27}(1,0),\mathbb F_7(3,1))$ is exactly
\[
 \mathbb P^1(\mathbb F_{27})\times\{\mathbb F_7(3,1)\}.
\]
The two separate implementations in Supplement, Section~\externalref{supplement}{app:computer} verify
the permutations and this entire finite orbit. Forward closure suffices
because inverses are positive powers of finite permutations. Orbit--stabilizer
therefore gives $28$, without a surjectivity assumption for the curve
group. Both $\Delta_0$ and $\Delta_H$ contain $-I$, so projectivization
preserves this index.
\end{proof}

\begin{theorem}[A connected principally polarized RM family]\label{thm:principalfamily}
The normalized modular embedding admits a connected finite cover $C'$
of its parametrizing orbifold and algebraic RM abelian schemes
$\mathcal A,\mathcal B$ of relative dimension six over $C'$, with an
$\OL$-linear isogeny
\[
 q:\mathcal A\longrightarrow\mathcal B,\qquad
 \deg(q)=453789,\qquad q^*\lambda_{\mathcal B}=\lambda_{\mathcal A}.
\]
Here $\lambda_{\mathcal A}$ is the original trace polarization and
$\lambda_{\mathcal B}$ is principal. The same connected cover contains
lifts of all three specified normalized vertices whose $\mathcal B$-fibers
are the common-lattice models of Proposition~\ref{prop:commonprincipal}.
Their specialized $\mathcal O_K$-actions and rational CM Hodge structures
are preserved, so the Hodge conjecture holds for these three fibers.

The moduli image lies in the principally polarized Hilbert--Blumenthal
component determined by the self-dual module $\Lambda_{\mathrm{common}}$.
No identification with Mostaed's named curve or global $K$-action is
included. The degree of a chosen auxiliary fine-level cover is not
asserted to be $28$.
\end{theorem}
\begin{proof}
Choose a neat finite-index subgroup $\Gamma_{\rm aux}\subset\Gamma_H$;
it exists by \cite[Thm.~3.6]{MilneModuli2011}, since $\Gamma_H$ is
arithmetic in $\operatorname{Res}_{L/\mathbb Q}\mathrm{SL}_2$.
Set $\Delta'=j^{-1}(\Gamma_{\rm aux})$ and
$C'=\mathbb H/\bar\Delta'$. The group $\bar\Delta'$ is torsion-free and
finite index in the cocompact quadrilateral group, so $C'$ is a connected
compact Riemann surface, hence a smooth projective algebraic curve.
Its map to the parametrizing orbifold is finite; the coarse map may
ramify at elliptic points. Since the neat group excludes $-I$, its
geometric degree is measured using projective group indices.

Over $\mathbb H$ use the fixed normalized embedding $F_t$ and the constant
lattices $\Lambda=\OL^2$ and $\Lambda'=\Lambda_{\mathrm{common}}$.
The transition action $\rho j$ descends both to integral local systems.
Their homology Hodge filtration $F^0$ is the kernel of the holomorphic map
\[
 (a,b)\longmapsto
 \bigl(\sigma_i(a)F_{t,i}(z)+\sigma_i(b)\bigr)_{i=1}^6.
\]
It has rank six and Hodge types $\{(-1,0),(0,-1)\}$; Griffiths
transversality is automatic for these types. The form $E_1$ is flat,
integral and positive by the Riemann-form identity already proved, and
is unimodular on $\Lambda'$. Thus both are polarizable integral
variations of Hodge structure. The discriminant local system
$\Lambda^\#/\Lambda$ has invariant subgroup $H$, and its inverse-image
lattice sheaf is exactly $\Lambda'$. This identifies the monodromy descent
of the common level.

Milne's equivalence for polarizable integral variations of these homology
types over smooth complex algebraic varieties
\cite[Thm.~7.13]{MilneModuli2011} algebraizes both families; the proof
explicitly includes nonprincipal polarizations. Full faithfulness
algebraizes the lattice inclusion, the $\OL$-actions and the
polarizations. The inclusion has index $453789$, and both forms are
$E_1$, giving the degree and pullback identity. Real multiplication is
self-adjoint and $\lambda_{\mathcal B}$ is principal.

Finally take the original three universal-cover parameters $z_A,z_B,z_C$.
Their images lie on the same connected curve $\mathbb H/\bar\Delta'$.
The constant marking $\Lambda'$ makes these fibers precisely the
three models in Proposition~\ref{prop:commonprincipal}. No choice of
components of a fiber product is involved. The specialized $K$-actions
need not be one global Hodge endomorphism. The RM family defines the
stated moduli map, and additional level gives a fine moduli interpretation.
\end{proof}
\begin{proposition}[The surviving order in the selected CM field]\label{prop:survivingorder}
For each of the three common principal fibers $B_X$, let
$\iota_X:M\hookrightarrow\operatorname{End}^0(B_X)$ be the selected CM
field embedding. Then
\[
 \operatorname{End}(B_X)\cap\iota_X(M)
 =\iota_X(R_K),\qquad
 R_K=\OL[\mathcal O_K]=\OL+\mathfrak p_3\OM.
\]
This order has index $27$ in $\OM$ and conductor $\mathfrak p_3\OM$.
The retained powers of the actual vertex automorphism are $0,7$ modulo
$14$ at $A$, and $0,7,14,21,28,35$ modulo $42$ at $B,C$.
The intersection assertion does not identify the full endomorphism
ring at the imprimitive $A$-fiber, or the generic endomorphism ring
of the family.
\end{proposition}
\begin{proof}
Put $z=\zeta_{42}$, so $\OM=\OL+\OL z$. Exact arithmetic gives
\[
 z^{14}=a+bz,\quad
 \begin{cases}
 a=-t^5+4t^3-3t-1,\\
 b=-t^5+t^4+6t^3-2t^2-8t-2.
 \end{cases}
\]
Both implementations verify $N_{L/\mathbb Q}(b)=-27$ and
$b\equiv0\pmod{\mathfrak p_3}$. Thus $(b)=\mathfrak p_3$, proving the
order formula and index. If $x=a_0+b_0z$ and both $x,xz$ lie in $R_K$,
their $z$ coefficients force $a_0,b_0\in\mathfrak p_3$; the reverse
inclusion is immediate. This proves the conductor assertion.

The intersection order $R'$ contains $R_K$ by the common-lattice theorem
and lies in $\OM$, since lattice endomorphisms are integral over $\mathbb Z$.
The quotient $\OM/R_K\simeq\OL/\mathfrak p_3$ is a simple $\OL$-module.
Consequently $R'$ is either $R_K$ or $\OM$. The latter would preserve the
full vertex automorphism, contradicted by the $\mathfrak p_3$ obstruction
proved above for all three vertices. Hence $R'=R_K$. Exhausting the
$14$ or $42$ powers on the two finite local subgroups gives the stated
surviving powers, with separate exact checks in both implementations.
\end{proof}

\begin{remark}[Imprimitive CM type at the $A$-vertex]
The stabilizer of $\Phi_A$ in $G=(\mathbb Z/42\mathbb Z)^\times$ is $\{1,29\}$.  Its fixed field is $\Ql(\zeta_7)$, and $\Phi_A$ is induced from the three-element CM type on that sextic CM field.  Thus the larger divisor algebra at the $A$-vertex is consistent with the fact that this CM type is imprimitive; in particular the associated sixfold is isogenous to the square of a CM threefold.
\end{remark}

